\documentclass[10pt,a4paper]{amsart}
\usepackage[margin=19mm]{geometry}
\usepackage{amsmath,amssymb,mathtools}
\usepackage[scr=rsfs,cal=euler]{mathalfa}
\usepackage{xfrac}
\usepackage[unicode,hidelinks,bookmarks=false]{hyperref}
\newcommand{\ie}{i.e.\ }
\newcommand{\cf}{cf.\ }
\newcommand{\resp}{respectively\ }
\newcommand{\Subsection}{\subsection}
\providecommand{\nobreakdash}{\nobreak\hbox{-}}
\setlength{\emergencystretch}{2em}
\multlinegap0pt
\usepackage{bm}
\usepackage{bbm}
\usepackage{dsfont}
\usepackage{xspace}
\usepackage{cancel}
\usepackage{mathtools}
\usepackage{amsthm}
\makeatletter
\@ifundefined{thm}{\newtheorem{thm}{Theorem}}{}
\@ifundefined{prop}{\newtheorem{prop}[thm]{Proposition}}{}
\makeatother
\title[A Generalization of a Classical Geometric Extremum Problem]{A Generalization of a Classical Geometric Extremum Problem}
\author{Petar Kenderov, Oleg Mushkarov,, Nikolai Nikolov}
\date{Source: arXiv:2506.07252v1 (CC BY 4.0)}
\begin{document}
\maketitle

\noindent\emph{Source and licence.} A Generalization of a Classical Geometric Extremum Problem. Version arXiv:2506.07252v1 (CC BY 4.0), distributed under the Creative Commons Attribution 4.0 International licence, as recorded in the arXiv licence metadata for this exact version and shown on the abstract page https://arxiv.org/abs/2506.07252v1. Full rights notice: licenses/arxiv-2506.07252-CC-BY-4.0.txt.\par\smallskip

\noindent\textit{Editorial scope.} Counted unit: the Thm whose source label is \texttt{thmmin} in the article (source0.tex lines 888--898, source label \texttt{thmmin}), delivered together with the article's own complete original proof of it (source0.tex lines 900--948, one \texttt{proof} environment of 49 lines). No mathematics is written, repaired or re-derived: statement and proof are the article's own text line for line. Required context retained uncounted: dim2prop (lines 391--402). Every definition, display and label of those lines is retained unchanged. Omitted from this package and not redistributed: 1--390, 403--887, 899--899, 949--1393. Nothing inside a retained excerpt is omitted or altered: every retained body is delivered exactly as the article's lines are written, so no omission receipt of any kind accompanies this package. Citations: the retained excerpts carry no citation command at all, so no bibliography entry of the article is transcribed and no reference is redistributed. Reference adaptation: none was needed and none was made; every reference inside the delivered unit resolves inside it, so no unresolved reference or citation is produced. Printed slips kept literal: no printed word, symbol, hypothesis or number of the counted statement or of its proof was altered. Layout and class: the article's own class is \texttt{amsart} with packages including amsmath, amssymb, graphicx, caption; the wrapper is the lane's pinned \texttt{amsart} editorial wrapper at 10pt on A4, which restates the theorem-like environments the retained text needs and adds the article's own macro definitions that the retained text needs (none), with extra wrapper packages (bm, bbm, dsfont, xspace, cancel, mathtools). The delivered numbering is the unit's own. Assets: no retained span contains a figure environment, a graphics-inclusion command, a PDF-page inclusion, a table or a TikZ picture; no image file is copied into this package, the package declares no asset dependency and no image file is redistributed with it. Licence: version arXiv:2506.07252v1 of this article is distributed under the Creative Commons Attribution 4.0 International licence, as recorded in the arXiv licence metadata for this exact version and shown on the abstract page https://arxiv.org/abs/2506.07252v1. A full rights notice is delivered at licenses/arxiv-2506.07252-CC-BY-4.0.txt. Characters: every retained excerpt is pure ASCII, so the delivered engine, XeLaTeX with its default Latin Modern text font, reports no missing character.
\begin{prop} \label{dim2prop}
Let $\mathcal C \subset \mathbb{R}^2$ be a compact convex set,  $O$ be an interior point of it, and
$\tilde{f}(\varphi) := |\tilde{A}(\varphi)\tilde{B}(\varphi)|$ where $[\tilde{A}(\varphi)\tilde{B}(\varphi)]=  l(\varphi)\cap \mathcal{C}$. Suppose $\varphi_*$ is a local minimum for $\tilde{f}(\varphi)$ and let $[A^*B^*]= l(\varphi_*) \cap \mathcal C$  Then:

(i) CPP holds for every pair of supporting for $\mathcal C$ lines
$l_1$ at $A^*$, $l_2$ at $B^*$ and the point $O \in [A^*B^*]$.

(ii) At each of the end points  $A^*$ and $B^*$ there exists only one supporting line for $\mathcal C$.

(iii) $\tilde{f}(\varphi)$ is differentiable at $\varphi_*$ and $\tilde{f}'(\varphi_*) = 0$.

\end{prop}

\begin{thm} \label{thmmin}
Let $O$ be an interior point of a closed convex set $\mathcal C
\subset \mathbb{R}^n,\, n \geq 2$, and let the segment $[A^*B^*]$
be a local minimizer for $F$. Then:

(i) There exists only one hyperplane $\Pi_1$   supporting $\mathcal C$   at  $A^*$ and only one hyperplane $\Pi_2$ supporting $\mathcal C$  at $B^*$;

(ii) The normal line to $\Pi_1$ at $A^*$ and the normal line to
$\Pi_2$ at $B^*$ intersect at a point belonging to the hyperplane
through $O$ which is orthogonal to the line $A^*B^*$.
\end{thm}

\begin{proof}
The particular case of this theorem when $n=2$ follows from
Proposition \ref{dim2prop}. Therefore, we will assume that $n
> 2$. Our first step is to show that  (ii) holds for arbitrary hyperplanes
$\Pi_1$ and $\Pi_2$ supporting $\mathcal C$ at $A^*$ and $B^*$ ,
respectively. Suppose first that $\Pi_1$ and $\Pi_2$ are not
parallel. Then there exist a linear subspace $L \subset
\mathbb{R}^n$ of co-dimension two and a point $T \in \Pi_1 \cap
\Pi_2 $ such that $\Pi_1 \cap \Pi_2 = \{T + h: h \in L\} = T + L$.
Take an arbitrary non-zero vector $h \in L $ and consider the
two-dimensional plane $\Pi(h)$ containing the lines $l_1(h)=\{A^*
+ th: t \in R\}$ and $l_2(h)=\{B^* + th: t \in R\}$. It is clear
that the parallel lines  $l_1(h)$ and $l_2(h)$ are supporting for
the set $\mathcal C \cap \Pi(h)$ in the plane $\Pi(h)$. Moreover,
$[A^*B^*]$ is a local minimizer for the the function $\tilde{f}$ corresponding to the set $\Pi(h) \cap \mathcal
C $. Then, as we have seen in the previous section, the segment $[A^*B^*]$ is a local minimizer for the function $f$ corresponding to the strip bounded by the lines $l_1(h)$ and $l_2(h)$. This means that $A^*B^*$ is perpendicular to  $l_1(h)$ and $l_2(h)$.
Since $h$ was an arbitrary vector from $L$ we
conclude that the line $A^*B^*$ is orthogonal to $L$ (and to
$T+L$).
 Hence there exists a two dimensional plane $\Pi$ which is orthogonal to $T+ L$ and contains the line $A^*B^*$. Put $l_1 = \Pi_1 \cap \Pi$ and $l_2 =\Pi_2 \cap \Pi$.
 The lines $l_1$ and $l_2$ are supporting in $\Pi$ the set $\mathcal C \cap \Pi$ at the points $A^*$ and $B^*$ , respectively, and $[A^*B^*]$ is a local minimizer for the set $\mathcal C \cap \Pi$  in $\Pi$.
 Proposition \ref{dim2prop} gives us three concurrent lines $p_0, p_1, p_2$ in $\Pi$ (with common point $P$)
 which are perpendicular , respectively, to the line $A^*B^*$ at $O$, to $l_1$ at $A^*$ and to $l_2$ at $B^*$.
 As $p_1$ lies in $\Pi$, it is orthogonal to $ T + L$. It is also perpendicular to $l_1$.
This implies that $p_1$ coincides with the normal line to $\Pi_1$ at $A^*$.
 Similarly, $p_2$ coincides with the normal line to $\Pi_2$ at $B^*$. Clearly, the line $p_0$, together with its point $P$ lies in the hyperplane through $O$ which is orthogonal to the line $A^*B^*$.
 Thus (ii)  holds for any nonparallel hyperplanes $\Pi_1$ and $\Pi_2$ supporting $\mathcal C$ at $A^*$ and $B^*$ , respectively.

 \smallskip


 Note that the line $p_1$, the normal to $\Pi_1$ at $A^*$, does not depend on the hyperplane $\Pi_2$ supporting $\mathcal C$ at $B^*$.
 Hence, the point $P$, as the intersection of $p_1$ and the hyperplane through $O$
 orthogonal to the line $A^*B^*$ is the same for every hyperplane $\Pi_2$ supporting $\mathcal C$ at $B^*$ and
 nonparallel to $\Pi_1$. This, in turn, means that any hyperplane $\Pi_2'$ supporting $\mathcal C$ at $B^*$ and
 not parallel to $\Pi_1$ is orthogonal to the line $PB^*$ and, therefore coincides with $\Pi_2$.

 In order to complete the proof, we have to discuss also the case where, along with $\Pi_2$,  there is another hyperplane $\Pi'$
 which supports $\mathcal C$ at $B^*$ and is parallel to $\Pi_1$.
 Suppose such a hyperplane $\Pi'$ exists. In this case, there is a linear subspace $L \subset \mathbb R^n$ of codimension one such that $\Pi_1 = A^* + L$ and $\Pi' = B^* + L$.
As above we can see that for every non zero vector $h \in L$ the
line $A^*B^*$ is perpendicular to the lines $l_1(h)=\{A^* + th: t
\in \mathbb{R}\}$ and $l_2(h)=\{B^* + th: t \in \mathbb{R}\}$.
This means that $A^*B^*$ is normal to both $\Pi_1$ and $\Pi'$ at
$A^*$ and $B^*$ , respectively. Thus $p_1 = A^*B^* $ and,
therefore, $P \in A^*B^*$. Since $PB^*$ is orthogonal to $\Pi_2$
we get $p_2 = A^*B^*$. This contradicts the assumption that
$\Pi_1$ and $\Pi_2$ are not parallel.
 \end{proof}

\end{document}
